% Compile with ceurart.cls from the official CEURART bundle:
% https://ceur-ws.org/Vol-XXX/CEURART.zip
\documentclass{ceurart}

\sloppy
\usepackage{booktabs}
\usepackage{hyperref}

\begin{document}

\copyrightyear{2026}
\copyrightclause{Copyright for this paper by its author. Use permitted under
  Creative Commons License Attribution 4.0 International (CC BY 4.0).}
\conference{Knowledge Engineering Automation Challenge 2026, co-located with
  WOP at ISWC 2026, Bari, Italy}

\title{KGGen+ for CQ2Term: Domain Graph Construction and Canonical Term Reassignment}

\author[1]{Philip Ganchev}[email=philip.ganchev@graphwise.ai]
\address[1]{Graphwise}

\begin{abstract}
We describe the Graphwise submission to the Knowledge Engineering Automation
Challenge 2026 CQ2Term sub-task. The system uses KGGen+ to construct a domain
knowledge graph from all competency questions, rather than predicting an
isolated term list from each question. Entities and relations are converted
into candidate ontology classes and properties while retaining question
provenance. String graph deduplication, lexical and morphological label
normalization, and conservative semantic clustering produce a canonical
domain vocabulary. A final semantic reassignment stage maps canonical terms
back to the questions that require them. The submitted configuration uses the
open 25.2B-parameter Gemma 4 26B A4B model and adds at most five terms of each
kind per question at a cosine threshold of 0.35. Across six repeated
public-data runs, reassignment leaves global term precision, recall, and F1
unchanged while improving mean question-conditioned coverage from 0.5901 to
0.7087. The submitted artifact is the run closest to the six-run mean. Source
code, exact commands, and version-pinned documentation are public.
\end{abstract}

\begin{keywords}
competency questions \sep ontology engineering \sep knowledge graphs \sep
large language models \sep term extraction \sep semantic clustering
\end{keywords}

\maketitle

\section{Introduction}

Competency questions (CQs) describe the information that an ontology should
represent and answer. The CQ2Term task isolates an early ontology-engineering
step: identifying the explicit classes and properties required by each CQ. A
useful system must recover the domain vocabulary while also assigning each
term to the correct question. The latter requirement is important because a
term may be present in the aggregate domain output but absent from the CQ that
needs it.

Our submission builds on KGGen+, an ontology-guided extension of KGGen
\cite{kggen}. Instead of treating each CQ as an independent zero-shot
term-prediction prompt, we first construct a graph over the complete CQ set of
a domain. This creates a shared vocabulary and preserves evidence linking
entities and relations to source questions. We then normalize and semantically
consolidate that vocabulary before reassigning canonical terms to individual
questions.

The submitted configuration was selected according to the CQ2Term
leaderboard's question-conditioned mean coverage criterion. It intentionally
does not include later experimental precision filters or predicate-induction
stages, because those improved global F1 while reducing question-conditioned
coverage.

\section{Method}

\subsection{Domain-wide graph construction}

For each of the six domains, all public competency questions are passed to
KGGen+ as separately identified input documents. The language model extracts
typed entities and subject--predicate--object relations from every CQ. Each
extracted item carries the identifier of the source question. Per-question
extraction is parallelized, after which the partial graphs are aggregated into
one domain graph.

The generation prompt treats common nouns, roles, states, events, and kinds as
candidate ontology concepts. Bracketed placeholders are interpreted as
references to classes rather than named individuals. Relations required to
represent the CQ are extracted as candidate properties. Generic fallback
predicates are discouraged.

The submitted graphs use deterministic string-based graph deduplication after
aggregation. Embedding-based graph deduplication is disabled because separate
public-data experiments showed that it could remove true terms and reduce
recall.

\subsection{Class and property projection}

CQ2Term classes are projected from entity surface forms. Properties are
projected from relation predicates. Provenance identifiers determine the
initial CQ-to-term assignments. Class and property vocabularies remain
separate throughout export.

The exporter applies Unicode and separator normalization, CamelCase splitting,
whitespace normalization, conservative class singularization, and first-verb
lemmatization for properties. Examples include \texttt{PlantParts} to
\texttt{plant part} and \texttt{uses platform} to \texttt{use platform}.

\subsection{Semantic clustering}

Normalized class labels and property labels are clustered separately with
\texttt{sentence-transformers/all-MiniLM-L6-v2}. Complete-link clustering at
cosine similarity 0.9 prevents a term from joining a cluster unless it is
sufficiently similar to every cluster member. Additional guards preserve
negation and prevent broader/narrower pairs such as \texttt{plant} and
\texttt{plant part} from being merged. The most frequent concise cluster
member becomes the canonical label.

Canonicalization is provenance-preserving. Runtime invariants assert that
every CQ that contributed an alias still receives the resulting canonical
term and that no unsupported CQ assignment is introduced.

\subsection{CQ-to-canonical-term reassignment}

Extraction and clustering preserve explicit provenance, but CQs containing
pronouns, ellipsis, or alternative wording may fail to receive a term already
discovered elsewhere in the domain. We therefore perform bounded semantic
reassignment.

For each CQ and term kind, the system embeds a query containing the CQ text and
its current assignments. Candidate terms are restricted to the canonical
vocabulary already extracted from that domain; reassignment cannot invent a
new term. Candidates with cosine similarity at least 0.35 are ranked, and at
most five additional classes and five additional properties are added to the
CQ.

This stage changes CQ-local attribution but not the set of unique domain
terms. It is therefore expected to improve CQ-conditioned coverage without
changing global term-level precision, recall, or F1.

\section{Experimental Setup}

The base model is
\texttt{cyankiwi/gemma-4-26B-A4B-it-AWQ-4bit}, an Apache-2.0-licensed
quantization of Gemma 4 26B A4B \cite{gemma}. The model has 25.2B total and
3.8B active parameters. It was served through vLLM under the OpenAI-compatible
model name \texttt{gemma4}.

Generation used temperature 0.7, a 4096-token completion limit, ten parallel
CQ documents, and six repeated runs. Each repeated run processed all six
domains: AWO, ODRL, SWO, VGO, Water, and Wine. Evaluation used the unmodified
CQ4OE evaluator distributed for the challenge \cite{cq4oe,challenge}.

The main ablation varied CQ reassignment:
\begin{itemize}
  \item no reassignment;
  \item threshold 0.55 with at most two additions (\texttt{t055-k2});
  \item threshold 0.45 with at most three additions (\texttt{t045-k3});
  \item threshold 0.35 with at most five additions (\texttt{t035-k5}).
\end{itemize}

The submitted \texttt{t035-k5} configuration had the highest public-data
CQ-Mean. Among its six repeated runs, run 03 had CQ-Mean 0.7079, closest to the
six-run mean of 0.7087, and was selected as a representative run rather than
by maximum score.

\section{Results}

Table~\ref{tab:results} reports six-run macro means for the submitted
configuration. Class and property F1 are global term-recovery metrics;
CQ-Any, CQ-Mean, and CQ-Full evaluate whether required terms occur under the
correct question.

\begin{table}[ht]
\caption{Six-run macro means for the submitted configuration.}
\label{tab:results}
\centering
\small
\begin{tabular}{lrrrrr}
\toprule
Domain & Class F1 & Property F1 & CQ-Any & CQ-Mean & CQ-Full \\
\midrule
AWO   & 0.9139 & 0.6667 & 1.0000 & 1.0000 & 1.0000 \\
ODRL  & 0.4722 & 0.4252 & 0.6316 & 0.3952 & 0.2632 \\
SWO   & 0.4538 & 0.5257 & 0.9679 & 0.6191 & 0.2564 \\
VGO   & 0.4159 & 0.4989 & 0.9545 & 0.8030 & 0.5682 \\
Water & 0.5613 & 0.7346 & 1.0000 & 0.6879 & 0.1083 \\
Wine  & 0.7871 & 0.3849 & 1.0000 & 0.7472 & 0.1000 \\
\midrule
Macro mean & \textbf{0.6007} & \textbf{0.5393} & \textbf{0.9257} &
\textbf{0.7087} & \textbf{0.3827} \\
\bottomrule
\end{tabular}
\end{table}

Without reassignment, the corresponding CQ-Any, CQ-Mean, and CQ-Full values
were 0.8813, 0.5901, and 0.2749. Global class and property metrics were
unchanged because reassignment only redistributes an existing canonical
vocabulary. Relative to no reassignment, \texttt{t035-k5} improved CQ-Any by
0.0444, CQ-Mean by 0.1187, and CQ-Full by 0.1077. Each coverage improvement
had the same direction in all six repeated runs; the exact two-sided sign-flip
$p$-value was 0.03125, the minimum non-zero value available with six pairs.

The threshold ablation was monotonic. CQ-Mean increased from 0.6631 for
\texttt{t055-k2}, to 0.6821 for \texttt{t045-k3}, and to 0.7087 for
\texttt{t035-k5}. CQ-Full similarly increased from 0.3310 to 0.3559 and
0.3827. This indicates that the more permissive policy recovered relevant
globally known terms without affecting the global vocabulary.

\section{Error Analysis and Limitations}

Performance differs substantially by domain. AWO has a small, lexically direct
vocabulary, while ODRL and SWO contain many ontology-specific property names
and implicit references. VGO achieves high question coverage but low class
precision because many plausible surface concepts are extracted beyond the
reference vocabulary. Water has the strongest property F1 but low CQ-Full,
indicating that properties are often discovered globally yet incompletely
assigned within individual questions.

Reassignment cannot recover a term that KGGen+ failed to extract anywhere in
the domain. It also cannot resolve a natural-language relation to an
ontology-specific identifier when no equivalent candidate exists. Conversely,
a low threshold may assign plausible but unnecessary terms to additional
questions. The cap limits this risk, but the current selection emphasizes
CQ-Mean rather than global F1.

Generation is stochastic. Temperature 0.7 was used to study repeated-run
variation, and byte-identical reproduction is not expected without reproducing
the complete model-serving environment and random-number state. The submitted
run was selected by closeness to the condition mean to reduce public-data
cherry-picking.

Finally, the method constructs a domain graph from the complete CQ set. This
is appropriate for shared ontology conceptualization but means that CQ outputs
are not independent. The approach intentionally exploits cross-question
domain context.

\section{Reproducibility}

KGGen+ source code is available under the MIT license at commit
\texttt{a7f0965}: \url{https://github.com/pgan002/kg-gen-plus/tree/a7f0965}.
Exact installation, vLLM serving, generation, export, evaluation, and packaging
commands are provided at
\url{https://github.com/pgan002/kg-gen-plus/blob/d5dd9df/docs/keac-2026-cq2term.md}.
The model is publicly available under Apache 2.0 at
\url{https://huggingface.co/cyankiwi/gemma-4-26B-A4B-it-AWQ-4bit}.
The challenge artifact is available in pull request 8:
\url{https://codeberg.org/ke-automation-challenge/challenge-catalog/pulls/8}.

\section{Conclusion}

KGGen+ frames CQ2Term as domain vocabulary construction followed by
provenance-aware question assignment. Canonical CQ reassignment substantially
improves question-conditioned coverage without altering global term-recovery
metrics. The strongest tested reassignment policy increases CQ-Mean by almost
twelve percentage points over the same extraction pipeline without
reassignment. Future work should improve ontology-specific predicate
canonicalization and conservative class selection while preserving the
coverage gains of domain-wide reassignment.

\begin{thebibliography}{9}

\bibitem{kggen}
B. Mo et al., ``KGGen: Extracting knowledge graphs from plain text with
language models.'' Source repository:
\url{https://github.com/stair-lab/kg-gen}.

\bibitem{cq4oe}
OEG-UPM, ``CQ4OE: A benchmark for assessing LLM-assisted ontology generation
from competency questions.''
\url{https://github.com/oeg-upm/cq4oe-benchmark}.

\bibitem{challenge}
Knowledge Engineering Automation Challenge 2026.
\url{https://ke-automation-challenge.codeberg.page/}.

\bibitem{sbert}
N. Reimers and I. Gurevych, ``Sentence-BERT: Sentence embeddings using Siamese
BERT-networks,'' in \emph{Proceedings of EMNLP-IJCNLP}, 2019.

\bibitem{gemma}
Google DeepMind, ``Gemma model family.''
\url{https://ai.google.dev/gemma}.

\end{thebibliography}

\end{document}
